% =====================================================================
%  Wrapper per compilare UNA lezione come PDF a sé stante.
%  Non si compila direttamente: ci pensa compila.ps1, che passa
%  \lezNum (numero della lezione) e \lezPath (percorso del sorgente):
%
%    pdflatex -jobname=lezione-02-utilita
%      "\def\lezNum{2}\def\lezPath{lezioni/02-utilita/lezione}% =====================================================================
%  Wrapper per compilare UNA lezione come PDF a sé stante.
%  Non si compila direttamente: ci pensa compila.ps1, che passa
%  \lezNum (numero della lezione) e \lezPath (percorso del sorgente):
%
%    pdflatex -jobname=lezione-02-utilita
%      "\def\lezNum{2}\def\lezPath{lezioni/02-utilita/lezione}% =====================================================================
%  Wrapper per compilare UNA lezione come PDF a sé stante.
%  Non si compila direttamente: ci pensa compila.ps1, che passa
%  \lezNum (numero della lezione) e \lezPath (percorso del sorgente):
%
%    pdflatex -jobname=lezione-02-utilita
%      "\def\lezNum{2}\def\lezPath{lezioni/02-utilita/lezione}% =====================================================================
%  Wrapper per compilare UNA lezione come PDF a sé stante.
%  Non si compila direttamente: ci pensa compila.ps1, che passa
%  \lezNum (numero della lezione) e \lezPath (percorso del sorgente):
%
%    pdflatex -jobname=lezione-02-utilita
%      "\def\lezNum{2}\def\lezPath{lezioni/02-utilita/lezione}\input{lezione-standalone}"
%
%  Lo stesso file lezione.tex finisce anche in main.tex, senza modifiche.
% =====================================================================
\documentclass[11pt, a4paper, oneside, openany]{book}
\input{preambolo}

\providecommand{\lezNum}{1}
\providecommand{\lezPath}{lezioni/01-introduzione/lezione}

% Il capitolo deve portare il numero della lezione anche da solo.
\setcounter{chapter}{\lezNum}
\addtocounter{chapter}{-1}

% L'intestazione va dentro al titolo del capitolo: \chapter apre comunque
% una pagina nuova, quindi metterla prima la lascerebbe da sola su pagina 1.
\newcommand{\intestazione}{%
  \begin{center}
    {\sffamily\bfseries\color{unical}\small UNIVERSITÀ DELLA CALABRIA\par}
    \vspace{0.15em}
    {\sffamily\color{grigio}\small Algorithmic Game Theory --- a.a.\ 2026/2027
     --- Sebastiano A.\ Piccolo\par}
    \vspace{0.15em}
    {\sffamily\color{grigio}\footnotesize Appunti in italiano, parte di
     \texttt{main.pdf} --- \today\par}
  \end{center}
  \vspace{-0.3em}
  {\color{blu}\rule{\linewidth}{1pt}}
  \vspace{1.2em}
  \par}

\titleformat{\chapter}[display]
  {\normalfont\sffamily\bfseries\color{blu}}
  {\intestazione\large\MakeUppercase{\chaptertitlename}\ \thechapter}
  {0.3em}
  {\Huge}
  [\vspace{0.3em}{\color{unical}\titlerule[1.2pt]}]

\begin{document}
\pagenumbering{arabic}
\input{\lezPath}
\end{document}
"
%
%  Lo stesso file lezione.tex finisce anche in main.tex, senza modifiche.
% =====================================================================
\documentclass[11pt, a4paper, oneside, openany]{book}
% =====================================================================
%  preambolo.tex — pacchetti, stile e ambienti comuni a tutto il documento
% =====================================================================

% ---------- Lingua e codifica ----------
\usepackage[T1]{fontenc}
\usepackage[utf8]{inputenc}
\usepackage[italian]{babel}

% ---------- Font ----------
% Times-like (newtx): tratto pieno e scuro, larghezza compatta.
% Anche la matematica è coerente col testo.
\usepackage{newtxtext}
\usepackage{newtxmath}
\usepackage[scaled=0.88]{beramono}   % monospace leggibile per il codice
\usepackage{microtype}

% ---------- Pagina ----------
\usepackage[a4paper, top=2.5cm, bottom=2.5cm, left=2.5cm, right=2.5cm, headheight=14pt]{geometry}
\setlength{\parindent}{0pt}
\setlength{\parskip}{0.55em}
\linespread{1.08}

% ---------- Colori ----------
\usepackage[dvipsnames]{xcolor}
\definecolor{unical}{HTML}{B00000}     % rosso delle slide
\definecolor{blu}{HTML}{0B3A91}        % blu delle slide di sezione
\definecolor{verde}{HTML}{1E7B34}
\definecolor{arancio}{HTML}{C0620A}
\definecolor{grigio}{HTML}{4A4A4A}
\definecolor{sfondo}{HTML}{F5F6F8}

% ---------- Matematica ----------
\usepackage{amsmath}
\usepackage{mathtools}

% ---------- Tabelle, liste, figure ----------
\usepackage{booktabs}
\usepackage{array}
\usepackage{multirow}
\usepackage{tabularx}
\usepackage{enumitem}
\setlist{itemsep=0.2em, topsep=0.3em}
\setlist[itemize,1]{label=\textcolor{unical}{\small$\blacktriangleright$}}
\usepackage{graphicx}
\usepackage{float}
\usepackage{caption}
\usepackage{needspace}   % \needspace{n\baselineskip}: va a capo pagina se manca spazio
\captionsetup{font=small, labelfont={bf,color=blu}}

\usepackage{tikz}
\usetikzlibrary{arrows.meta, positioning, trees, calc, shapes.geometric}

% ---------- Titoli ----------
\usepackage{titlesec}
% Capitolo = una lezione
\titleformat{\chapter}[display]
  {\normalfont\sffamily\bfseries\color{blu}}
  {\large\MakeUppercase{\chaptertitlename}\ \thechapter}
  {0.3em}
  {\Huge}
  [\vspace{0.3em}{\color{unical}\titlerule[1.2pt]}]
\titlespacing*{\chapter}{0pt}{0pt}{1.5em}
\titleformat{\section}
  {\normalfont\sffamily\Large\bfseries\color{blu}}{\thesection}{0.8em}{}
\titleformat{\subsection}
  {\normalfont\sffamily\large\bfseries\color{black}}{\thesubsection}{0.8em}{}
\titleformat{\subsubsection}
  {\normalfont\sffamily\normalsize\bfseries\color{grigio}}{\thesubsubsection}{0.8em}{}
\addto\captionsitalian{\renewcommand{\chaptername}{Lezione}}

% ---------- Intestazioni ----------
\usepackage{fancyhdr}
\pagestyle{fancy}
\fancyhf{}
\fancyhead[L]{\small\sffamily\color{grigio}Algorithmic Game Theory}
\fancyhead[R]{\small\sffamily\color{grigio}\leftmark}
\fancyfoot[C]{\small\sffamily\thepage}
\renewcommand{\headrulewidth}{0.4pt}
\renewcommand{\chaptermark}[1]{\markboth{Lezione \thechapter\ -- #1}{}}
\fancypagestyle{plain}{\fancyhf{}\fancyfoot[C]{\small\sffamily\thepage}\renewcommand{\headrulewidth}{0pt}}

% ---------- Box (definizioni, esempi, ...) ----------
\usepackage[most]{tcolorbox}
\tcbset{
  appunti/.style={
    enhanced, breakable, sharp corners,
    boxrule=0pt, leftrule=3pt,
    left=8pt, right=8pt, top=5pt, bottom=5pt,
    before skip=0.8em, after skip=0.8em,
    fonttitle=\sffamily\bfseries, coltitle=black,
    attach title to upper, after title={.\ },
  }
}
\newtcolorbox[auto counter, number within=chapter]{definizione}[1][]{
  appunti, colback=blu!5, colframe=blu,
  title={\textcolor{blu}{Definizione~\thetcbcounter}\ifx\relax#1\relax\else\ (#1)\fi}}
\newtcolorbox[use counter from=definizione]{teorema}[1][]{
  appunti, colback=unical!5, colframe=unical,
  title={\textcolor{unical}{Teorema~\thetcbcounter}\ifx\relax#1\relax\else\ (#1)\fi}}
\newtcolorbox{esempio}[1][]{
  appunti, colback=verde!5, colframe=verde,
  title={\textcolor{verde}{Esempio}\ifx\relax#1\relax\else\ (#1)\fi}}
\newtcolorbox{osservazione}[1][]{
  appunti, colback=sfondo, colframe=grigio,
  title={\textcolor{grigio}{Osservazione}\ifx\relax#1\relax\else\ (#1)\fi}}
% Punti delicati, errori tipici, imprecisioni nelle slide
\newtcolorbox{attenzione}[1][]{
  appunti, colback=arancio!7, colframe=arancio,
  title={\textcolor{arancio}{Attenzione}\ifx\relax#1\relax\else\ (#1)\fi}}
% Contenuti aggiunti negli appunti, non presenti nelle slide
\newtcolorbox{extra}[1][]{
  appunti, colback=white, colframe=grigio!60, leftrule=1.5pt,
  borderline west={0.4pt}{4pt}{grigio!60},
  title={\textcolor{grigio}{Approfondimento}\ifx\relax#1\relax\else\ (#1)\fi}}
% Riepilogo di fine lezione
\newtcolorbox{riepilogo}{
  enhanced, breakable, sharp corners, colback=blu!4, colframe=blu, boxrule=0.8pt,
  fonttitle=\sffamily\bfseries, title={Riepilogo della lezione}}

% ---------- Teoria vs Laboratorio ----------
% \parteteoria e \partelab aprono le due metà di una lezione con un'etichetta colorata.
\newcommand{\etichetta}[2]{%
  \tikz[baseline=(t.base)]\node[fill=#1, text=white, font=\sffamily\bfseries\small,
    inner xsep=6pt, inner ysep=3pt](t){#2};}
\newcommand{\parteteoria}{%
  \par\medskip\noindent\etichetta{blu}{TEORIA}\par\smallskip}
\newcommand{\partelab}[1]{%
  \clearpage
  \section{Laboratorio: #1}
  \noindent\etichetta{unical}{LAB}\par\smallskip}
\newcommand{\fonte}[1]{\hfill{\small\sffamily\color{grigio}[#1]}}

% ---------- Codice ----------
\usepackage{listings}
\lstdefinestyle{python}{
  language=Python,
  basicstyle=\ttfamily\small,
  keywordstyle=\color{blu}\bfseries,
  commentstyle=\color{verde}\itshape,
  stringstyle=\color{unical},
  numberstyle=\tiny\color{grigio}, numbers=left, numbersep=8pt,
  backgroundcolor=\color{sfondo},
  frame=leftline, framerule=2pt, rulecolor=\color{blu},
  xleftmargin=1.8em, framexleftmargin=1.8em,
  showstringspaces=false, breaklines=true, tabsize=4,
  columns=fullflexible, keepspaces=true,
  morekeywords={True,False,None},
  literate={à}{{\`a}}1 {è}{{\`e}}1 {é}{{\'e}}1 {ì}{{\`i}}1 {ò}{{\`o}}1 {ù}{{\`u}}1
           {−}{{-}}1 {∞}{{$\infty$}}1 {×}{{$\times$}}1,
}
\lstdefinestyle{pseudo}{
  style=python, numbers=none, mathescape=true,
  morekeywords={def,if,else,for,each,return,break,in,and,or,not},
}
\lstset{style=python}

% ---------- Tris (tic-tac-toe) ----------
% \tris{X,O,e, e,X,e, O,e,e}: 9 celle riga per riga, e = vuota.
\newcommand{\trisX}{\textbf{X}}
\newcommand{\trisO}{\textcolor{unical}{\textbf{O}}}
\newcommand{\trise}{}
\newcommand{\tris}[2][0.42]{%
  \begin{tikzpicture}[scale=#1, baseline=(current bounding box.center)]
    \draw[thick] (1,0)--(1,3) (2,0)--(2,3) (0,1)--(3,1) (0,2)--(3,2);
    \foreach \c [count=\k from 0] in {#2} {
      \pgfmathtruncatemacro\riga{\k/3}
      \pgfmathtruncatemacro\colonna{mod(\k,3)}
      \node[font=\sffamily] at (\colonna+0.5, 2.5-\riga) {\csname tris\c\endcsname};
    }
  \end{tikzpicture}}

% ---------- Utilità ----------
\newcommand{\en}[1]{\textit{#1}}            % termine inglese originale
\newcommand{\kw}[1]{\textbf{#1}}            % parola chiave
\newcommand{\R}{\mathbb{R}}
\newcommand{\N}{\mathbb{N}}
\newcommand{\Ev}{\mathbb{E}}
% Relazioni di preferenza
\newcommand{\pref}{\succsim}                % x pref y: x e' almeno buono quanto y
\newcommand{\npref}{\not\succsim}
\newcommand{\spref}{\succ}                  % preferenza stretta
\newcommand{\indif}{\sim}                   % indifferenza
\newcommand{\lott}{\mathcal{L}}             % insieme delle lotterie semplici
% Teoria dei giochi
\newcommand{\BR}{\mathrm{BR}}               % best response
\newcommand{\uno}{\mathbf{1}}               % vettore di tutti 1
\DeclareMathOperator{\supp}{supp}           % supporto di una strategia mista
\DeclareMathOperator{\core}{Core}
% Tabella dei payoff di un gioco a due giocatori: celle "a, b"
\newcommand{\pay}[2]{$#1,\,#2$}
\DeclareMathOperator*{\argmax}{arg\,max}
\DeclareMathOperator*{\argmin}{arg\,min}

% ---------- Link (per ultimo) ----------
\usepackage[hidelinks]{hyperref}
\hypersetup{pdftitle={Algorithmic Game Theory -- Appunti}, pdflang={it}}
\usepackage{bookmark}


\providecommand{\lezNum}{1}
\providecommand{\lezPath}{lezioni/01-introduzione/lezione}

% Il capitolo deve portare il numero della lezione anche da solo.
\setcounter{chapter}{\lezNum}
\addtocounter{chapter}{-1}

% L'intestazione va dentro al titolo del capitolo: \chapter apre comunque
% una pagina nuova, quindi metterla prima la lascerebbe da sola su pagina 1.
\newcommand{\intestazione}{%
  \begin{center}
    {\sffamily\bfseries\color{unical}\small UNIVERSITÀ DELLA CALABRIA\par}
    \vspace{0.15em}
    {\sffamily\color{grigio}\small Algorithmic Game Theory --- a.a.\ 2026/2027
     --- Sebastiano A.\ Piccolo\par}
    \vspace{0.15em}
    {\sffamily\color{grigio}\footnotesize Appunti in italiano, parte di
     \texttt{main.pdf} --- \today\par}
  \end{center}
  \vspace{-0.3em}
  {\color{blu}\rule{\linewidth}{1pt}}
  \vspace{1.2em}
  \par}

\titleformat{\chapter}[display]
  {\normalfont\sffamily\bfseries\color{blu}}
  {\intestazione\large\MakeUppercase{\chaptertitlename}\ \thechapter}
  {0.3em}
  {\Huge}
  [\vspace{0.3em}{\color{unical}\titlerule[1.2pt]}]

\begin{document}
\pagenumbering{arabic}
\input{\lezPath}
\end{document}
"
%
%  Lo stesso file lezione.tex finisce anche in main.tex, senza modifiche.
% =====================================================================
\documentclass[11pt, a4paper, oneside, openany]{book}
% =====================================================================
%  preambolo.tex — pacchetti, stile e ambienti comuni a tutto il documento
% =====================================================================

% ---------- Lingua e codifica ----------
\usepackage[T1]{fontenc}
\usepackage[utf8]{inputenc}
\usepackage[italian]{babel}

% ---------- Font ----------
% Times-like (newtx): tratto pieno e scuro, larghezza compatta.
% Anche la matematica è coerente col testo.
\usepackage{newtxtext}
\usepackage{newtxmath}
\usepackage[scaled=0.88]{beramono}   % monospace leggibile per il codice
\usepackage{microtype}

% ---------- Pagina ----------
\usepackage[a4paper, top=2.5cm, bottom=2.5cm, left=2.5cm, right=2.5cm, headheight=14pt]{geometry}
\setlength{\parindent}{0pt}
\setlength{\parskip}{0.55em}
\linespread{1.08}

% ---------- Colori ----------
\usepackage[dvipsnames]{xcolor}
\definecolor{unical}{HTML}{B00000}     % rosso delle slide
\definecolor{blu}{HTML}{0B3A91}        % blu delle slide di sezione
\definecolor{verde}{HTML}{1E7B34}
\definecolor{arancio}{HTML}{C0620A}
\definecolor{grigio}{HTML}{4A4A4A}
\definecolor{sfondo}{HTML}{F5F6F8}

% ---------- Matematica ----------
\usepackage{amsmath}
\usepackage{mathtools}

% ---------- Tabelle, liste, figure ----------
\usepackage{booktabs}
\usepackage{array}
\usepackage{multirow}
\usepackage{tabularx}
\usepackage{enumitem}
\setlist{itemsep=0.2em, topsep=0.3em}
\setlist[itemize,1]{label=\textcolor{unical}{\small$\blacktriangleright$}}
\usepackage{graphicx}
\usepackage{float}
\usepackage{caption}
\usepackage{needspace}   % \needspace{n\baselineskip}: va a capo pagina se manca spazio
\captionsetup{font=small, labelfont={bf,color=blu}}

\usepackage{tikz}
\usetikzlibrary{arrows.meta, positioning, trees, calc, shapes.geometric}

% ---------- Titoli ----------
\usepackage{titlesec}
% Capitolo = una lezione
\titleformat{\chapter}[display]
  {\normalfont\sffamily\bfseries\color{blu}}
  {\large\MakeUppercase{\chaptertitlename}\ \thechapter}
  {0.3em}
  {\Huge}
  [\vspace{0.3em}{\color{unical}\titlerule[1.2pt]}]
\titlespacing*{\chapter}{0pt}{0pt}{1.5em}
\titleformat{\section}
  {\normalfont\sffamily\Large\bfseries\color{blu}}{\thesection}{0.8em}{}
\titleformat{\subsection}
  {\normalfont\sffamily\large\bfseries\color{black}}{\thesubsection}{0.8em}{}
\titleformat{\subsubsection}
  {\normalfont\sffamily\normalsize\bfseries\color{grigio}}{\thesubsubsection}{0.8em}{}
\addto\captionsitalian{\renewcommand{\chaptername}{Lezione}}

% ---------- Intestazioni ----------
\usepackage{fancyhdr}
\pagestyle{fancy}
\fancyhf{}
\fancyhead[L]{\small\sffamily\color{grigio}Algorithmic Game Theory}
\fancyhead[R]{\small\sffamily\color{grigio}\leftmark}
\fancyfoot[C]{\small\sffamily\thepage}
\renewcommand{\headrulewidth}{0.4pt}
\renewcommand{\chaptermark}[1]{\markboth{Lezione \thechapter\ -- #1}{}}
\fancypagestyle{plain}{\fancyhf{}\fancyfoot[C]{\small\sffamily\thepage}\renewcommand{\headrulewidth}{0pt}}

% ---------- Box (definizioni, esempi, ...) ----------
\usepackage[most]{tcolorbox}
\tcbset{
  appunti/.style={
    enhanced, breakable, sharp corners,
    boxrule=0pt, leftrule=3pt,
    left=8pt, right=8pt, top=5pt, bottom=5pt,
    before skip=0.8em, after skip=0.8em,
    fonttitle=\sffamily\bfseries, coltitle=black,
    attach title to upper, after title={.\ },
  }
}
\newtcolorbox[auto counter, number within=chapter]{definizione}[1][]{
  appunti, colback=blu!5, colframe=blu,
  title={\textcolor{blu}{Definizione~\thetcbcounter}\ifx\relax#1\relax\else\ (#1)\fi}}
\newtcolorbox[use counter from=definizione]{teorema}[1][]{
  appunti, colback=unical!5, colframe=unical,
  title={\textcolor{unical}{Teorema~\thetcbcounter}\ifx\relax#1\relax\else\ (#1)\fi}}
\newtcolorbox{esempio}[1][]{
  appunti, colback=verde!5, colframe=verde,
  title={\textcolor{verde}{Esempio}\ifx\relax#1\relax\else\ (#1)\fi}}
\newtcolorbox{osservazione}[1][]{
  appunti, colback=sfondo, colframe=grigio,
  title={\textcolor{grigio}{Osservazione}\ifx\relax#1\relax\else\ (#1)\fi}}
% Punti delicati, errori tipici, imprecisioni nelle slide
\newtcolorbox{attenzione}[1][]{
  appunti, colback=arancio!7, colframe=arancio,
  title={\textcolor{arancio}{Attenzione}\ifx\relax#1\relax\else\ (#1)\fi}}
% Contenuti aggiunti negli appunti, non presenti nelle slide
\newtcolorbox{extra}[1][]{
  appunti, colback=white, colframe=grigio!60, leftrule=1.5pt,
  borderline west={0.4pt}{4pt}{grigio!60},
  title={\textcolor{grigio}{Approfondimento}\ifx\relax#1\relax\else\ (#1)\fi}}
% Riepilogo di fine lezione
\newtcolorbox{riepilogo}{
  enhanced, breakable, sharp corners, colback=blu!4, colframe=blu, boxrule=0.8pt,
  fonttitle=\sffamily\bfseries, title={Riepilogo della lezione}}

% ---------- Teoria vs Laboratorio ----------
% \parteteoria e \partelab aprono le due metà di una lezione con un'etichetta colorata.
\newcommand{\etichetta}[2]{%
  \tikz[baseline=(t.base)]\node[fill=#1, text=white, font=\sffamily\bfseries\small,
    inner xsep=6pt, inner ysep=3pt](t){#2};}
\newcommand{\parteteoria}{%
  \par\medskip\noindent\etichetta{blu}{TEORIA}\par\smallskip}
\newcommand{\partelab}[1]{%
  \clearpage
  \section{Laboratorio: #1}
  \noindent\etichetta{unical}{LAB}\par\smallskip}
\newcommand{\fonte}[1]{\hfill{\small\sffamily\color{grigio}[#1]}}

% ---------- Codice ----------
\usepackage{listings}
\lstdefinestyle{python}{
  language=Python,
  basicstyle=\ttfamily\small,
  keywordstyle=\color{blu}\bfseries,
  commentstyle=\color{verde}\itshape,
  stringstyle=\color{unical},
  numberstyle=\tiny\color{grigio}, numbers=left, numbersep=8pt,
  backgroundcolor=\color{sfondo},
  frame=leftline, framerule=2pt, rulecolor=\color{blu},
  xleftmargin=1.8em, framexleftmargin=1.8em,
  showstringspaces=false, breaklines=true, tabsize=4,
  columns=fullflexible, keepspaces=true,
  morekeywords={True,False,None},
  literate={à}{{\`a}}1 {è}{{\`e}}1 {é}{{\'e}}1 {ì}{{\`i}}1 {ò}{{\`o}}1 {ù}{{\`u}}1
           {−}{{-}}1 {∞}{{$\infty$}}1 {×}{{$\times$}}1,
}
\lstdefinestyle{pseudo}{
  style=python, numbers=none, mathescape=true,
  morekeywords={def,if,else,for,each,return,break,in,and,or,not},
}
\lstset{style=python}

% ---------- Tris (tic-tac-toe) ----------
% \tris{X,O,e, e,X,e, O,e,e}: 9 celle riga per riga, e = vuota.
\newcommand{\trisX}{\textbf{X}}
\newcommand{\trisO}{\textcolor{unical}{\textbf{O}}}
\newcommand{\trise}{}
\newcommand{\tris}[2][0.42]{%
  \begin{tikzpicture}[scale=#1, baseline=(current bounding box.center)]
    \draw[thick] (1,0)--(1,3) (2,0)--(2,3) (0,1)--(3,1) (0,2)--(3,2);
    \foreach \c [count=\k from 0] in {#2} {
      \pgfmathtruncatemacro\riga{\k/3}
      \pgfmathtruncatemacro\colonna{mod(\k,3)}
      \node[font=\sffamily] at (\colonna+0.5, 2.5-\riga) {\csname tris\c\endcsname};
    }
  \end{tikzpicture}}

% ---------- Utilità ----------
\newcommand{\en}[1]{\textit{#1}}            % termine inglese originale
\newcommand{\kw}[1]{\textbf{#1}}            % parola chiave
\newcommand{\R}{\mathbb{R}}
\newcommand{\N}{\mathbb{N}}
\newcommand{\Ev}{\mathbb{E}}
% Relazioni di preferenza
\newcommand{\pref}{\succsim}                % x pref y: x e' almeno buono quanto y
\newcommand{\npref}{\not\succsim}
\newcommand{\spref}{\succ}                  % preferenza stretta
\newcommand{\indif}{\sim}                   % indifferenza
\newcommand{\lott}{\mathcal{L}}             % insieme delle lotterie semplici
% Teoria dei giochi
\newcommand{\BR}{\mathrm{BR}}               % best response
\newcommand{\uno}{\mathbf{1}}               % vettore di tutti 1
\DeclareMathOperator{\supp}{supp}           % supporto di una strategia mista
\DeclareMathOperator{\core}{Core}
% Tabella dei payoff di un gioco a due giocatori: celle "a, b"
\newcommand{\pay}[2]{$#1,\,#2$}
\DeclareMathOperator*{\argmax}{arg\,max}
\DeclareMathOperator*{\argmin}{arg\,min}

% ---------- Link (per ultimo) ----------
\usepackage[hidelinks]{hyperref}
\hypersetup{pdftitle={Algorithmic Game Theory -- Appunti}, pdflang={it}}
\usepackage{bookmark}


\providecommand{\lezNum}{1}
\providecommand{\lezPath}{lezioni/01-introduzione/lezione}

% Il capitolo deve portare il numero della lezione anche da solo.
\setcounter{chapter}{\lezNum}
\addtocounter{chapter}{-1}

% L'intestazione va dentro al titolo del capitolo: \chapter apre comunque
% una pagina nuova, quindi metterla prima la lascerebbe da sola su pagina 1.
\newcommand{\intestazione}{%
  \begin{center}
    {\sffamily\bfseries\color{unical}\small UNIVERSITÀ DELLA CALABRIA\par}
    \vspace{0.15em}
    {\sffamily\color{grigio}\small Algorithmic Game Theory --- a.a.\ 2026/2027
     --- Sebastiano A.\ Piccolo\par}
    \vspace{0.15em}
    {\sffamily\color{grigio}\footnotesize Appunti in italiano, parte di
     \texttt{main.pdf} --- \today\par}
  \end{center}
  \vspace{-0.3em}
  {\color{blu}\rule{\linewidth}{1pt}}
  \vspace{1.2em}
  \par}

\titleformat{\chapter}[display]
  {\normalfont\sffamily\bfseries\color{blu}}
  {\intestazione\large\MakeUppercase{\chaptertitlename}\ \thechapter}
  {0.3em}
  {\Huge}
  [\vspace{0.3em}{\color{unical}\titlerule[1.2pt]}]

\begin{document}
\pagenumbering{arabic}
\input{\lezPath}
\end{document}
"
%
%  Lo stesso file lezione.tex finisce anche in main.tex, senza modifiche.
% =====================================================================
\documentclass[11pt, a4paper, oneside, openany]{book}
% =====================================================================
%  preambolo.tex — pacchetti, stile e ambienti comuni a tutto il documento
% =====================================================================

% ---------- Lingua e codifica ----------
\usepackage[T1]{fontenc}
\usepackage[utf8]{inputenc}
\usepackage[italian]{babel}

% ---------- Font ----------
% Times-like (newtx): tratto pieno e scuro, larghezza compatta.
% Anche la matematica è coerente col testo.
\usepackage{newtxtext}
\usepackage{newtxmath}
\usepackage[scaled=0.88]{beramono}   % monospace leggibile per il codice
\usepackage{microtype}

% ---------- Pagina ----------
\usepackage[a4paper, top=2.5cm, bottom=2.5cm, left=2.5cm, right=2.5cm, headheight=14pt]{geometry}
\setlength{\parindent}{0pt}
\setlength{\parskip}{0.55em}
\linespread{1.08}

% ---------- Colori ----------
\usepackage[dvipsnames]{xcolor}
\definecolor{unical}{HTML}{B00000}     % rosso delle slide
\definecolor{blu}{HTML}{0B3A91}        % blu delle slide di sezione
\definecolor{verde}{HTML}{1E7B34}
\definecolor{arancio}{HTML}{C0620A}
\definecolor{grigio}{HTML}{4A4A4A}
\definecolor{sfondo}{HTML}{F5F6F8}

% ---------- Matematica ----------
\usepackage{amsmath}
\usepackage{mathtools}

% ---------- Tabelle, liste, figure ----------
\usepackage{booktabs}
\usepackage{array}
\usepackage{multirow}
\usepackage{tabularx}
\usepackage{enumitem}
\setlist{itemsep=0.2em, topsep=0.3em}
\setlist[itemize,1]{label=\textcolor{unical}{\small$\blacktriangleright$}}
\usepackage{graphicx}
\usepackage{float}
\usepackage{caption}
\usepackage{needspace}   % \needspace{n\baselineskip}: va a capo pagina se manca spazio
\captionsetup{font=small, labelfont={bf,color=blu}}

\usepackage{tikz}
\usetikzlibrary{arrows.meta, positioning, trees, calc, shapes.geometric}

% ---------- Titoli ----------
\usepackage{titlesec}
% Capitolo = una lezione
\titleformat{\chapter}[display]
  {\normalfont\sffamily\bfseries\color{blu}}
  {\large\MakeUppercase{\chaptertitlename}\ \thechapter}
  {0.3em}
  {\Huge}
  [\vspace{0.3em}{\color{unical}\titlerule[1.2pt]}]
\titlespacing*{\chapter}{0pt}{0pt}{1.5em}
\titleformat{\section}
  {\normalfont\sffamily\Large\bfseries\color{blu}}{\thesection}{0.8em}{}
\titleformat{\subsection}
  {\normalfont\sffamily\large\bfseries\color{black}}{\thesubsection}{0.8em}{}
\titleformat{\subsubsection}
  {\normalfont\sffamily\normalsize\bfseries\color{grigio}}{\thesubsubsection}{0.8em}{}
\addto\captionsitalian{\renewcommand{\chaptername}{Lezione}}

% ---------- Intestazioni ----------
\usepackage{fancyhdr}
\pagestyle{fancy}
\fancyhf{}
\fancyhead[L]{\small\sffamily\color{grigio}Algorithmic Game Theory}
\fancyhead[R]{\small\sffamily\color{grigio}\leftmark}
\fancyfoot[C]{\small\sffamily\thepage}
\renewcommand{\headrulewidth}{0.4pt}
\renewcommand{\chaptermark}[1]{\markboth{Lezione \thechapter\ -- #1}{}}
\fancypagestyle{plain}{\fancyhf{}\fancyfoot[C]{\small\sffamily\thepage}\renewcommand{\headrulewidth}{0pt}}

% ---------- Box (definizioni, esempi, ...) ----------
\usepackage[most]{tcolorbox}
\tcbset{
  appunti/.style={
    enhanced, breakable, sharp corners,
    boxrule=0pt, leftrule=3pt,
    left=8pt, right=8pt, top=5pt, bottom=5pt,
    before skip=0.8em, after skip=0.8em,
    fonttitle=\sffamily\bfseries, coltitle=black,
    attach title to upper, after title={.\ },
  }
}
\newtcolorbox[auto counter, number within=chapter]{definizione}[1][]{
  appunti, colback=blu!5, colframe=blu,
  title={\textcolor{blu}{Definizione~\thetcbcounter}\ifx\relax#1\relax\else\ (#1)\fi}}
\newtcolorbox[use counter from=definizione]{teorema}[1][]{
  appunti, colback=unical!5, colframe=unical,
  title={\textcolor{unical}{Teorema~\thetcbcounter}\ifx\relax#1\relax\else\ (#1)\fi}}
\newtcolorbox{esempio}[1][]{
  appunti, colback=verde!5, colframe=verde,
  title={\textcolor{verde}{Esempio}\ifx\relax#1\relax\else\ (#1)\fi}}
\newtcolorbox{osservazione}[1][]{
  appunti, colback=sfondo, colframe=grigio,
  title={\textcolor{grigio}{Osservazione}\ifx\relax#1\relax\else\ (#1)\fi}}
% Punti delicati, errori tipici, imprecisioni nelle slide
\newtcolorbox{attenzione}[1][]{
  appunti, colback=arancio!7, colframe=arancio,
  title={\textcolor{arancio}{Attenzione}\ifx\relax#1\relax\else\ (#1)\fi}}
% Contenuti aggiunti negli appunti, non presenti nelle slide
\newtcolorbox{extra}[1][]{
  appunti, colback=white, colframe=grigio!60, leftrule=1.5pt,
  borderline west={0.4pt}{4pt}{grigio!60},
  title={\textcolor{grigio}{Approfondimento}\ifx\relax#1\relax\else\ (#1)\fi}}
% Riepilogo di fine lezione
\newtcolorbox{riepilogo}{
  enhanced, breakable, sharp corners, colback=blu!4, colframe=blu, boxrule=0.8pt,
  fonttitle=\sffamily\bfseries, title={Riepilogo della lezione}}

% ---------- Teoria vs Laboratorio ----------
% \parteteoria e \partelab aprono le due metà di una lezione con un'etichetta colorata.
\newcommand{\etichetta}[2]{%
  \tikz[baseline=(t.base)]\node[fill=#1, text=white, font=\sffamily\bfseries\small,
    inner xsep=6pt, inner ysep=3pt](t){#2};}
\newcommand{\parteteoria}{%
  \par\medskip\noindent\etichetta{blu}{TEORIA}\par\smallskip}
\newcommand{\partelab}[1]{%
  \clearpage
  \section{Laboratorio: #1}
  \noindent\etichetta{unical}{LAB}\par\smallskip}
\newcommand{\fonte}[1]{\hfill{\small\sffamily\color{grigio}[#1]}}

% ---------- Codice ----------
\usepackage{listings}
\lstdefinestyle{python}{
  language=Python,
  basicstyle=\ttfamily\small,
  keywordstyle=\color{blu}\bfseries,
  commentstyle=\color{verde}\itshape,
  stringstyle=\color{unical},
  numberstyle=\tiny\color{grigio}, numbers=left, numbersep=8pt,
  backgroundcolor=\color{sfondo},
  frame=leftline, framerule=2pt, rulecolor=\color{blu},
  xleftmargin=1.8em, framexleftmargin=1.8em,
  showstringspaces=false, breaklines=true, tabsize=4,
  columns=fullflexible, keepspaces=true,
  morekeywords={True,False,None},
  literate={à}{{\`a}}1 {è}{{\`e}}1 {é}{{\'e}}1 {ì}{{\`i}}1 {ò}{{\`o}}1 {ù}{{\`u}}1
           {−}{{-}}1 {∞}{{$\infty$}}1 {×}{{$\times$}}1,
}
\lstdefinestyle{pseudo}{
  style=python, numbers=none, mathescape=true,
  morekeywords={def,if,else,for,each,return,break,in,and,or,not},
}
\lstset{style=python}

% ---------- Tris (tic-tac-toe) ----------
% \tris{X,O,e, e,X,e, O,e,e}: 9 celle riga per riga, e = vuota.
\newcommand{\trisX}{\textbf{X}}
\newcommand{\trisO}{\textcolor{unical}{\textbf{O}}}
\newcommand{\trise}{}
\newcommand{\tris}[2][0.42]{%
  \begin{tikzpicture}[scale=#1, baseline=(current bounding box.center)]
    \draw[thick] (1,0)--(1,3) (2,0)--(2,3) (0,1)--(3,1) (0,2)--(3,2);
    \foreach \c [count=\k from 0] in {#2} {
      \pgfmathtruncatemacro\riga{\k/3}
      \pgfmathtruncatemacro\colonna{mod(\k,3)}
      \node[font=\sffamily] at (\colonna+0.5, 2.5-\riga) {\csname tris\c\endcsname};
    }
  \end{tikzpicture}}

% ---------- Utilità ----------
\newcommand{\en}[1]{\textit{#1}}            % termine inglese originale
\newcommand{\kw}[1]{\textbf{#1}}            % parola chiave
\newcommand{\R}{\mathbb{R}}
\newcommand{\N}{\mathbb{N}}
\newcommand{\Ev}{\mathbb{E}}
% Relazioni di preferenza
\newcommand{\pref}{\succsim}                % x pref y: x e' almeno buono quanto y
\newcommand{\npref}{\not\succsim}
\newcommand{\spref}{\succ}                  % preferenza stretta
\newcommand{\indif}{\sim}                   % indifferenza
\newcommand{\lott}{\mathcal{L}}             % insieme delle lotterie semplici
% Teoria dei giochi
\newcommand{\BR}{\mathrm{BR}}               % best response
\newcommand{\uno}{\mathbf{1}}               % vettore di tutti 1
\DeclareMathOperator{\supp}{supp}           % supporto di una strategia mista
\DeclareMathOperator{\core}{Core}
% Tabella dei payoff di un gioco a due giocatori: celle "a, b"
\newcommand{\pay}[2]{$#1,\,#2$}
\DeclareMathOperator*{\argmax}{arg\,max}
\DeclareMathOperator*{\argmin}{arg\,min}

% ---------- Link (per ultimo) ----------
\usepackage[hidelinks]{hyperref}
\hypersetup{pdftitle={Algorithmic Game Theory -- Appunti}, pdflang={it}}
\usepackage{bookmark}


\providecommand{\lezNum}{1}
\providecommand{\lezPath}{lezioni/01-introduzione/lezione}

% Il capitolo deve portare il numero della lezione anche da solo.
\setcounter{chapter}{\lezNum}
\addtocounter{chapter}{-1}

% L'intestazione va dentro al titolo del capitolo: \chapter apre comunque
% una pagina nuova, quindi metterla prima la lascerebbe da sola su pagina 1.
\newcommand{\intestazione}{%
  \begin{center}
    {\sffamily\bfseries\color{unical}\small UNIVERSITÀ DELLA CALABRIA\par}
    \vspace{0.15em}
    {\sffamily\color{grigio}\small Algorithmic Game Theory --- a.a.\ 2026/2027
     --- Sebastiano A.\ Piccolo\par}
    \vspace{0.15em}
    {\sffamily\color{grigio}\footnotesize Appunti in italiano, parte di
     \texttt{main.pdf} --- \today\par}
  \end{center}
  \vspace{-0.3em}
  {\color{blu}\rule{\linewidth}{1pt}}
  \vspace{1.2em}
  \par}

\titleformat{\chapter}[display]
  {\normalfont\sffamily\bfseries\color{blu}}
  {\intestazione\large\MakeUppercase{\chaptertitlename}\ \thechapter}
  {0.3em}
  {\Huge}
  [\vspace{0.3em}{\color{unical}\titlerule[1.2pt]}]

\begin{document}
\pagenumbering{arabic}
\input{\lezPath}
\end{document}
